% BHU PYQ -- Transcribed & Corrected
% Paper: BHUVA-13: Ayurveda -- Traditional Health Care & Management Principles (Sem II Mid Term)
% Exam: Under Graduate Semester II Mid-Term Examination, 2024-25 / 2025-26 (NEP)
% Category: Value Added Course (VAC)

\documentclass[11pt,a4paper]{article}
\usepackage[a4paper,top=2.5cm,bottom=2.5cm,left=2.8cm,right=2.8cm,headheight=14pt]{geometry}
\usepackage{amsmath,amssymb}
\usepackage[T1]{fontenc}
\usepackage[utf8]{inputenc}
\usepackage{lmodern,microtype,fancyhdr,enumitem,hyperref,lastpage,parskip}

\hypersetup{
    colorlinks=true,
    urlcolor=black,
    linkcolor=black,
    pdftitle={BHUVA13: Ayurveda Sem II Mid-Term -- BHU 2025-26 (NEP VAC)}
}

\pagestyle{fancy}
\fancyhf{}
\renewcommand{\headrulewidth}{0.4pt}
\renewcommand{\footrulewidth}{0.4pt}
\fancyhead[L]{\small   \textit{Banaras Hindu University}}
\fancyhead[R]{\small   \textit{BHUVA-13: Ayurveda (Sem II Mid Term VAC)}}
\fancyfoot[C]{\small Page  \thepage\ of \pageref{LastPage}}


\newlist{parts}{enumerate}{2}
\setlist[parts,1]{label=(\Alph*),leftmargin=*,itemsep=3pt,topsep=2pt}

\newcommand{\pts}[1]{\hfill   \textit{[#1]}}


\begin{document}


\begin{center}
  {\large\bfseries BANARAS HINDU UNIVERSITY}\\[2pt]
  {\normalsize Under Graduate --- Semester II Mid-Term Examination, 2024-25 / 2025-26 (NEP)}\\[6pt]
  \rule{0.8\linewidth}{0.5pt}\\[6pt]
  {\large\bfseries Value Added Course (VAC)}\\[4pt]
  {\large\bfseries Paper Code: BHUVA-13 : Ayurveda --- Traditional Health Care \& Management Principles (Sem II Mid Term)}\\[4pt]
  {\normalsize Time: 1 Hour\hfill Maximum Marks: 30}
\end{center}

\rule{\linewidth}{0.4pt}

\smallskip



\noindent   \textit{Instructions: Answer all 30 Multiple Choice Questions. Each question carries equal marks (1 mark each).}

\bigskip




\noindent   \textbf{Question 1.} How many purification procedures are included in Shodhana Chikitsa (Panchakarma)?\pts{1}

\begin{parts}
  \item Four
  \item Five (Vamana, Virechana, Basti, Nasya, Raktamokshana)
  \item Six
  \item Seven
\end{parts}

\medskip




\noindent   \textbf{Question 2.} Which stage of Shodhana Karma involves Deepana-Pachana, Snehana, and Swedana?\pts{1}

\begin{parts}
  \item Pradhana Karma
  \item Poorva Karma (Preparatory phase)
  \item Paschat Karma
  \item Upakrama Karma
\end{parts}

\medskip




\noindent   \textbf{Question 3.} Which procedure is NOT included in Panchashodhana?\pts{1}

\begin{parts}
  \item Vamana
  \item Virechana
  \item Raktamokshana
  \item Anuvasana Basti
\end{parts}

\medskip




\noindent   \textbf{Question 4.} What are the four essential components of Ayu as defined in Ayurveda?\pts{1}

\begin{parts}
  \item Body, Soul, Knowledge, Virtue
  \item Body, Mind, Soul, and Sense Organs (   \textit{Sharira, Indriya, Satva, Atma})
  \item Health, Strength, Vitality, Immunity
  \item Consciousness, Soul, Body, Intelligence
\end{parts}

\medskip




\noindent   \textbf{Question 5.} Which branch of Ayurveda deals with pediatrics and childcare?\pts{1}

\begin{parts}
  \item Kaumarbhritya
  \item Kaya Chikitsa
  \item Shalya Tantra
  \item Agada Tantra
\end{parts}

\medskip




\noindent   \textbf{Question 6.} Dhatus in Ayurveda refer to:\pts{1}

\begin{parts}
  \item Structural and functional tissues of the body
  \item Excretory waste products
  \item Toxic substances
  \item Environmental elements
\end{parts}

\medskip




\noindent   \textbf{Question 7.} Which Dhatu is responsible for blood circulation and life vitality (   \textit{Jeevana})?\pts{1}

\begin{parts}
  \item Rakta Dhatu
  \item Rasa Dhatu
  \item Mamsa Dhatu
  \item Asthi Dhatu
\end{parts}

\medskip




\noindent   \textbf{Question 8.} Which Agni is responsible for primary digestion in the gastrointestinal tract?\pts{1}

\begin{parts}
  \item Jatharagni
  \item Bhutagni
  \item Dhatvagni
  \item Mandagni
\end{parts}

\medskip




\noindent   \textbf{Question 9.} Which of the following is considered a Sadvritta (righteous conduct) practice?\pts{1}

\begin{parts}
  \item Speaking the truth and respecting elders
  \item Overeating
  \item Late-night sleeping
  \item Suppressing natural urges
\end{parts}

\medskip




\noindent   \textbf{Question 10.} Which Ayurvedic classical text deals mainly with Kaumarbhritya (pediatrics)?\pts{1}

\begin{parts}
  \item Kashyap Samhita
  \item Charaka Samhita
  \item Sushruta Samhita
  \item Astanga Sangraha
\end{parts}

\medskip




\noindent   \textbf{Question 11.} What are the principles of Ahara (Diet) according to Ayurveda?\pts{1}

\begin{parts}
  \item Food should be balanced, fresh, and nourishing
  \item It should be consumed mindfully with proper time and temperature
  \item It should suit one's Prakriti (individual constitution)
  \item All of the above
\end{parts}

\medskip




\noindent   \textbf{Question 12.} What is the main purpose of Karnavedhana Samskara (ear piercing)?\pts{1}

\begin{parts}
  \item For protection and health benefits
  \item For decoration
  \item Both (A) and (B)
  \item None of the above
\end{parts}

\medskip




\noindent   \textbf{Question 13.} In Ayurveda, female reproductive physiology and menstruation is understood by the term:\pts{1}

\begin{parts}
  \item Dosha
  \item Dhatu
  \item Artava
  \item Shukra
\end{parts}

\medskip




\noindent   \textbf{Question 14.} Factors essential for maintaining reproductive and sexual health as per Ayurveda are:\pts{1}

\begin{parts}
  \item Maintaining mental, emotional, and physical balance
  \item Regular intake of spicy foods
  \item Avoiding salt in diet
  \item Fasting frequently
\end{parts}

\medskip




\noindent   \textbf{Question 15.} Which therapy is used for rejuvenation, longevity, and memory/vitality enhancement in Ayurveda?\pts{1}

\begin{parts}
  \item Rasayana Therapy
  \item Shodhana
  \item Shamana
  \item Langhana
\end{parts}

\medskip




\noindent   \textbf{Question 16.} Why is tongue cleaning (   \textit{Jihva Nirlekhana}) an important Ayurvedic practice?\pts{1}

\begin{parts}
  \item It freshens breath, removes toxic coating (   \textit{Ama}), and enhances digestion
  \item It prevents bad dreams
  \item It strengthens teeth
  \item None of the above
\end{parts}

\medskip




\noindent   \textbf{Question 17.} The meaning of '   \textit{Kaya}' in Kaya Chikitsa is:\pts{1}

\begin{parts}
  \item Agni / Body / Metabolism
  \item Shalya Karma
  \item Clothes
  \item Aged person
\end{parts}

\medskip




\noindent   \textbf{Question 18.} Which para-surgical procedure is a well-known Ayurvedic method for treating anal fistula (   \textit{Bhagandara})?\pts{1}

\begin{parts}
  \item Ksharasutra
  \item Agnikarma
  \item Raktamokshana
  \item Chedana
\end{parts}

\medskip




\noindent   \textbf{Question 19.} Which therapy is NOT a part of Shalakya Tantra (ENT \& Ophthalmology) treatment?\pts{1}

\begin{parts}
  \item Ksharasutra
  \item Nasya
  \item Karna Poorana
  \item Netra Tarpana
\end{parts}

\medskip




\noindent   \textbf{Question 20.} In Ayurveda, which daily practice is recommended for preventing diseases affecting the eyes?\pts{1}

\begin{parts}
  \item Netra Prakshalana (Eye washing)
  \item Nasya
  \item Karnapoorana
  \item Virechana
\end{parts}

\medskip




\noindent   \textbf{Question 21.} Why is tongue cleaning an important part of Dinacharya?\pts{1}

\begin{parts}
  \item Removes toxic coating (   \textit{Ama}) and enhances taste perception
  \item Strengthens teeth
  \item Clears throat
  \item None of the above
\end{parts}

\medskip




\noindent   \textbf{Question 22.} Direct perception using sense organs is called:\pts{1}

\begin{parts}
  \item Pratyaksha Pramana
  \item Anumana Pramana
  \item Yukti
  \item Aaptopadesha
\end{parts}

\medskip




\noindent   \textbf{Question 23.} Inference based on cause and effect reasoning is known as:\pts{1}

\begin{parts}
  \item Anumana Pramana
  \item Pratyaksha Pramana
  \item Aaptopadesha
  \item Upamana
\end{parts}

\medskip




\noindent   \textbf{Question 24.} Which Dosha is responsible for movement, enthusiasm, and respiratory activities in the body?\pts{1}

\begin{parts}
  \item Vata Dosha
  \item Pitta Dosha
  \item Kapha Dosha
  \item None of the above
\end{parts}

\medskip




\noindent   \textbf{Question 25.} Which Dosha regulates digestion, body temperature, and metabolic heat?\pts{1}

\begin{parts}
  \item Pitta Dosha
  \item Vata Dosha
  \item Kapha Dosha
  \item All of the above
\end{parts}

\medskip




\noindent   \textbf{Question 26.} Samhitas (Ayurvedic classical texts) are considered as:\pts{1}

\begin{parts}
  \item Aaptopadesha (Authoritative testimonial knowledge)
  \item Pratyaksha
  \item Anumana
  \item Yukti
\end{parts}

\medskip




\noindent   \textbf{Question 27.} Pratyaksha, Anumana, and Aaptopadesha are called:\pts{1}

\begin{parts}
  \item Pramanas (Means of valid knowledge)
  \item Pariksha
  \item Both A \& B
  \item None of the above
\end{parts}

\medskip




\noindent   \textbf{Question 28.} Which dimension is NOT included by WHO in the definition of Health?\pts{1}

\begin{parts}
  \item Spiritual / Financial dimension
  \item Physical dimension
  \item Mental dimension
  \item Social dimension
\end{parts}

\medskip




\noindent   \textbf{Question 29.} Direct acknowledgment using sense organs is called:\pts{1}

\begin{parts}
  \item Pratyaksha
  \item Anumana
  \item Yukti
  \item Aaptopadesha
\end{parts}

\medskip




\noindent   \textbf{Question 30.} Nasya therapy involves the administration of medicated oils or powders through the:\pts{1}

\begin{parts}
  \item Nasal passages (   \textit{Nasa Hi Shirso Dvaram})
  \item Oral cavity
  \item Ears
  \item Eyes
\end{parts}

\vfill
\rule{\linewidth}{0.4pt}

\begin{center}\small
\href{https://bhu-exam.vercel.app/}{Click here to visit our website}\\[4pt]
\href{https://me-aryan.vercel.app/}{Click here to visit our contributor}
\end{center}

\end{document}
